\section{A finitely presented affine group}\label{sec:affine}

We now use the algebraic data in Theorem~\ref{thm:algebra-package} to
construct a finitely presented group of translations and elementary
operations.  The translations form the additive group of $B^6$.
We first give their finite presentation relative to a Steinberg group,
and then construct the endomorphism used in the next two sections.

\subsection{The coefficient ring and the translation module}

Put
\[
 R=B\otimes_k B^{\op},
 \qquad (a\otimes b^{\op})\cdot m=amb\quad(m\in B).
\]
Thus $r=\sum_\nu a_\nu\otimes b_\nu^{\op}$ acts on $B$ by the
operation $m\mapsto\sum_\nu a_\nu m b_\nu$. The opposite multiplication
in the second factor makes composition of these operations agree with
multiplication in $R$, so $B$ is a left $R$-module.
The ring $R$ is finitely presented
as a unital associative ring: take a finite algebra presentation of
$B$, a second copy with products reversed, the relations commuting
every generator of the first copy with every generator of the second,
and $2\cdot1=0$.  The universal property of the tensor product
identifies the resulting ring with $R$.
The module $B$ is generated by $1_B$; we next show that the relations
on this module generator also have a finite description.

\begin{lemma}\label{lem:cyclic-module}
Let $Y$ be a finite algebra generating set for $B$.  For $b\in B$ put
\[
 d_b=b\otimes1-1\otimes b^{\op},
 \qquad J=\sum_{y\in Y}R d_y.
\]
Evaluation at $1_B$ induces an isomorphism of left $R$-modules
$R/J\cong B$.
\end{lemma}

\begin{proof}
Evaluation is a surjective $R$-module homomorphism, and every $d_y$
annihilates $1_B$.  The identity
\begin{equation}\label{eq:module-product}
 d_{bc}=(b\otimes1)d_c+(1\otimes c^{\op})d_b
\end{equation}
shows, together with additivity and $d_1=0$, that $d_b\in J$ for all
$b\in B$.  Indeed, $B$ is generated by $Y$ over the prime field $k$.
Consequently
\[
 a\otimes b^{\op}\equiv ab\otimes1\pmod J
 \qquad(a,b\in B).
\]
Every element of $R$ is therefore congruent modulo $J$ to its
evaluation tensored with $1$.  The additive map
$b\mapsto(b\otimes1)+J$ is an inverse to the induced evaluation map.
Thus the kernel is exactly the stated left ideal.
\end{proof}

For a unital ring $R$ and $n\geq3$, the Steinberg group $\St_n(R)$
has generators $x_{ab}(r)$, where $1\leq a\ne b\leq n$ and $r\in R$,
with relations
\begin{align}
 x_{ab}(r)x_{ab}(s)&=x_{ab}(r+s),\label{eq:steinberg-add}\\
 [x_{ab}(r),x_{bc}(s)]&=x_{ac}(rs)
       &&(a,b,c\text{ distinct}),\label{eq:steinberg-product}\\
 [x_{ab}(r),x_{cd}(s)]&=1
       &&(a\ne d,\ b\ne c).\label{eq:steinberg-commute}
\end{align}
Let $e_{ab}$ denote the matrix unit.  The map
$x_{ab}(r)\mapsto I+r e_{ab}$ gives a surjection
\[
 \pi_n:\St_n(R)\longrightarrow\E_n(R),
\]
where $\E_n(R)$ is the subgroup of $\mathrm{GL}_n(R)$ generated by
these elementary matrices.  Krsti\'c and McCool prove that
$\St_n(R)$ is finitely presented whenever $R$ is a finitely presented
unital associative ring and $n\geq4$
\cite[Theorem~3]{KrsticMcCool1999}.  Their theorem permits
noncommutative rings, and applies to the ring constructed above.

\subsection{A finite relative presentation}

The next lemma supplies the additional relations needed for the
translations.  Finite presentation of the module and of the linear
group alone would not establish its conclusion; we prove the
semidirect-product presentation explicitly.

\begin{lemma}\label{lem:affine-presentation}
Let $R$ be a finitely generated unital associative ring, let
$J=\sum_{d\in D}Rd$ for a finite subset $D\subset R$, and let $n\geq4$.
If $\St_n(R)$ is finitely presented, then
\[
 (R/J)^n\rtimes\St_n(R)
\]
is finitely presented, where the action on columns is induced by
$\pi_n$.  More precisely, choose a finite ring generating set
$F\subset R$ containing $1$.  A finite presentation is obtained
from a finite presentation of $\St_n(R)$ by adjoining
$p_1,\ldots,p_n$ and the relations
\begin{align}
 [p_a,p_b]&=1,\label{eq:affine-p}\\
 [x_{ab}(s),p_c]&=1
     &&(s\in F,\ c\ne b),\label{eq:affine-central}\\
 [x_{ab}(1),p_b]&=p_a,\label{eq:affine-unit}\\
 [x_{ab}(d),p_b]&=1
     &&(d\in D).\label{eq:affine-ideal}
\end{align}
Here every root has $a\ne b$, and each root occurring in these
finitely many relations is represented by a word in the chosen
generators of $\St_n(R)$.
\end{lemma}

\begin{proof}
Let $Q$ be the group given by this presentation. In the intended
semidirect product, $p_a$ translates by $1+J$ in row $a$, and
$[x_{ab}(r),p_b]$ translates by $r+J$ in that row. We recover these
translations and their elementary action from the finite relations,
working inside $Q$ until the final inverse-homomorphism argument.

\paragraph{Coefficients.}
Relation~\eqref{eq:affine-central} holds for every $s\in R$.
To see this, let $C$ be the set of coefficients for which it holds
for every permitted triple of indices.  By
\eqref{eq:steinberg-add}, $C$ is an additive subgroup containing
$F$.  If $r,s\in C$, choose $h\notin\{a,b,c\}$.  There is such an
index because $n\geq4$.  Then
\[
 x_{ab}(rs)=[x_{ah}(r),x_{hb}(s)],
\]
and both factors commute with $p_c$.  Thus $C$ is closed under
multiplication, so $C=R$.

\paragraph{Independence of a column.}
For $a\ne b$, define
\[
 w_a^b(r)=[x_{ab}(r),p_b].
\]
If $h$ differs from $a,b$, then $w_a^b(r)$ commutes with $p_h$:
both $x_{ab}(r)$ and $p_b$ do.  Put
\[
 X=x_{ah}(r),\qquad Y=x_{hb}(1),\qquad Z=x_{ab}(r).
\]
The Steinberg relation $[X,Y]=Z$ gives $XY=ZYX$.
Conjugating $p_b$ by these equal products yields
\begin{align*}
 (XY)p_b(XY)^{-1}&=w_a^h(r)p_hp_b,\\
 (ZYX)p_b(ZYX)^{-1}&=p_hw_a^b(r)p_b.
\end{align*}
Here $X$ fixes $p_b$, $Y$ sends $p_b$ to $p_hp_b$, and $Z$ fixes
$p_h$.  Cancelling $p_b$ and using the preceding commutation with
$p_h$ proves $w_a^h(r)=w_a^b(r)$.  We write this common value as
$w_a(r)$.

\paragraph{Commutations.}
The element $w_a(r)$ commutes with every $p_c$: choose its defining
column $b\notin\{a,c\}$, so that both defining factors commute with
$p_c$.  It also commutes with $x_{cd}(s)$ whenever $d\ne a$.
For this, choose the defining column $b\notin\{a,c,d\}$.
Then $x_{ab}(r)$ commutes with $x_{cd}(s)$ by
\eqref{eq:steinberg-commute}, and $p_b$ commutes with $x_{cd}(s)$
because $b\ne d$.

These facts imply that all the $w$'s commute.  Given $w_a(r)$ and
$w_c(s)$, choose $b\notin\{a,c\}$ and write
$w_c(s)=[x_{cb}(s),p_b]$.  The element $w_a(r)$ commutes with both
factors, hence with their commutator.  This argument includes $a=c$.

\paragraph{Additivity.}
Using $[uv,p]=u[v,p]u^{-1}[u,p]$ and any $b\ne a$, we obtain
\begin{align*}
 w_a(r+s)
  &=[x_{ab}(r)x_{ab}(s),p_b]\\
  &=x_{ab}(r)w_a(s)x_{ab}(r)^{-1}w_a(r)\\
  &=w_a(r)w_a(s).
\end{align*}
The last equality uses the commutations just proved.  Thus
$w_a:(R,+)\to Q$ is a homomorphism.  Moreover,
$w_a(1)=p_a$ by \eqref{eq:affine-unit}.  In particular, any additive
torsion of $R$ is already imposed on these translations; in
characteristic two, $p_a^2=w_a(1+1)=1$.

\paragraph{Conjugation.}
It remains to determine the action of roots whose column is $a$.
For $c\ne a$, choose $b\notin\{a,c\}$ and put
\[
 X=x_{ca}(s),\qquad Y=x_{ab}(r),\qquad Z=x_{cb}(sr).
\]
The element $X$ fixes $p_b$, and $XYX^{-1}=ZY$.
Consequently
\begin{align}
 Xw_a(r)X^{-1}
    &=[ZY,p_b]
      =Zw_a(r)Z^{-1}w_c(sr)
      =w_a(r)w_c(sr).\label{eq:affine-action}
\end{align}
The last equality holds because the column $b$ of $Z$ differs
from $a$.  Together with the commutations already proved,
\eqref{eq:affine-action} is precisely the elementary action on
an additive column over the left module $R$.

\paragraph{The module relations.}
Relation~\eqref{eq:affine-ideal} says $w_a(d)=1$ for every $a$ and
$d\in D$.  To kill a left multiple in a prescribed row $c$, choose
$a\ne c$ and conjugate $w_a(d)=1$ by $x_{ca}(s)$.
Equation~\eqref{eq:affine-action} gives $w_c(sd)=1$.
Additivity now proves that every $w_c$ vanishes on $J$.

\paragraph{Inverse homomorphisms.}
Let $P_0=(R/J)^n\rtimes\St_n(R)$.  Sending $p_a$ to translation by
$1+J$ in row $a$ and using the identity on $\St_n(R)$ satisfies
all the defining relations of $Q$, and hence gives a homomorphism
$f:Q\to P_0$.  Conversely, the commuting additive maps
\[
 r+J\longmapsto w_a(r)
\]
define a homomorphism $(R/J)^n\to Q$.
Equation~\eqref{eq:affine-action} and the root commutations give
the required equivariance, so this map and the canonical map
$\St_n(R)\to Q$ induce $g:P_0\to Q$.
The composite $gf$ fixes $\St_n(R)$ and each
$p_a=w_a(1)$, so it is the identity on $Q$.
The composite $fg$ fixes $\St_n(R)$ and every translation:
in $P_0$, the commutator of $x_{ab}(r)$ with translation by $1+J$
in row $b$ is translation by $r+J$ in row $a$.
Thus $f$ and $g$ are inverse isomorphisms.
\end{proof}

\begin{proposition}\label{prop:affine}
For the algebra $B$ in Theorem~\ref{thm:algebra-package}, the group
\[
 \Gamma_0=B^6\rtimes\St_6(R),\qquad R=B\otimes_kB^{\op},
\]
is finitely presented.
\end{proposition}

\begin{proof}
The ring $R$ is finitely presented, so the Krsti\'c--McCool theorem
applies to $\St_6(R)$.  Lemma~\ref{lem:cyclic-module} identifies $B$
with a quotient of $R$ by a finitely generated left ideal.
Lemma~\ref{lem:affine-presentation}, with $n=6$, gives the result.
\end{proof}

\subsection{The transition endomorphism}

Write $\St=\St_6(R)$, $\E=\E_6(R)$, and
$\pi=\pi_6:\St\to\E$.  Recall the injective algebra map
$\phi:B\to B$ from Theorem~\ref{thm:algebra-package}, and set
$\psi=\phi\otimes\phi^{\op}:R\to R$.
This map is injective because tensoring vector spaces over $k$
preserves injections.  Both maps may be nonunital.

\begin{lemma}\label{lem:transition}
The formulas
\[
 \theta(x_{ab}(r))=x_{ab}(\psi(r)),\qquad
 \Theta(m,h)=(\phi^{(6)}(m),\theta(h))
\]
define endomorphisms of $\St$ and $\Gamma_0$, respectively.
Here $\phi^{(6)}$ acts componentwise on $B^6$.
The induced map on elementary matrices is the injective
homomorphism
\begin{equation}\label{eq:matrix-transition}
 F:\E\longrightarrow\E,
 \qquad F(g)=I+\psi^{(6)}(g-I),
\end{equation}
where $\psi^{(6)}$ means entrywise application on $6\times6$
matrices.  Thus $\pi\theta=F\pi$.
\end{lemma}

\begin{proof}
Additivity and multiplicativity of $\psi$ preserve all the
Steinberg relations, proving that $\theta$ is a homomorphism.
For an elementary tensor $r=a\otimes b^{\op}$ and $m\in B$,
\[
 \phi(r\cdot m)=\phi(amb)
   =\phi(a)\phi(m)\phi(b)
   =\psi(r)\cdot\phi(m).
\]
By additivity this holds for every $r\in R$.  Consequently, for
each elementary root,
\[
 \phi^{(6)}\bigl((I+r e_{ab})m\bigr)
   =(I+\psi(r)e_{ab})\phi^{(6)}(m).
\]
The roots generate $\St$, so this equivariance proves that
$\Theta$ is an endomorphism of the semidirect product.

For the matrix assertion, write $g=I+A$ and $h=I+C$.
Entrywise application of $\psi$ is additive and multiplicative
on matrix rings, even though it does not preserve their identity.
Since $gh-I=A+C+AC$, we have
\begin{align*}
 F(gh)
  &=I+\psi^{(6)}(A)+\psi^{(6)}(C)
       +\psi^{(6)}(A)\psi^{(6)}(C)\\
  &=F(g)F(h).
\end{align*}
Also $F(I)=I$, and the same calculation shows
$F(g^{-1})=F(g)^{-1}$.  On elementary matrices,
$F(I+r e_{ab})=I+\psi(r)e_{ab}$, proving both that $F$ maps
$\E$ into $\E$ and that $\pi\theta=F\pi$.
Finally, $F(g)=I$ forces every entry of $g-I$ to vanish because
$\psi$ is injective.  Thus $F$ is injective.
\end{proof}

We have obtained a finite presentation and compatible transition
maps.  The map $\Theta$ need not be injective: the remaining
obstruction lies in the kernel of $\pi$.  The next section proves
that $\theta$ annihilates this entire kernel.
